\section{Data and Corpus Construction}

\subsection{Source context}

The corpus context follows the post-junta parliamentary archive setting described in earlier work: written parliamentary questions (1976--1977), predominantly in formal Katharevousa. Upstream OCR and document processing create usable text but leave substantial reconstruction demands for syntax-oriented NLP.

\subsection{Reconstruction and preprocessing}

Before annotation, we apply reconstruction-focused preprocessing to repair OCR-induced segmentation and format errors, especially:
\begin{itemize}[leftmargin=1.25em]
  \item line-break hyphenation and split-word joins,
  \item malformed punctuation around sentence boundaries,
  \item long enumerative sentences that break naive splitting,
  \item parser-breaking malformed JSON from annotation calls.
\end{itemize}

\subsection{LLM-assisted gold generation and freeze}

Annotation is generated through schema-constrained GPT batch processing with retry and resume logic. A deterministic freeze step merges batch and retry outputs into one deduplicated final snapshot:
\begin{itemize}[leftmargin=1.25em]
  \item final sentences: 1,697,
  \item source breakdown: 1,565 batch + 132 retry replacements.
\end{itemize}

The frozen file and manifest ensure exact reruns:
\begin{itemize}[leftmargin=1.25em]
  \item \texttt{data/processed/final\_gold/gold\_final.conllu}
  \item \texttt{data/processed/final\_gold/snapshot\_manifest.json}
\end{itemize}
